\paragraph{Scope and provenance.}
This is an independent proof and continuation of Section 5 of
\texttt{Tau\_Gamma\_Joint\_Schur\_Return/NOTE.tex} in the supplied
14 September 2026 package. Its determinant and finite bounds are correct.
The expansion proved below strengthens that section while retaining its
actual measure, real evaluation nodes, polynomial degree, and all four
endpoints. It evaluates the free boundary contribution. The original
relation-dependent Schur complement remains a separate summand in the
exact signed return, and is not assigned the value of this contribution.

\section{The actual source and its polynomial evaluation covariance}
Let $l\ge1$, $k=4l+1$, $c=k/2$, $\vartheta=1/2$, and
\[
 \sigma(y)=\frac{|\Gamma(1/4+iy/2)|^2}{2\pi},\qquad
 c_\sigma=\sqrt{2\pi},\qquad
 \alpha_j=c+2j+\tfrac12\quad(0\le j<l).
\]
The polynomial space $\mathcal P_M=\{P\in\C[S]:\deg P\le M\}$ carries
the original inner product
\[
 \langle P,Q\rangle_\sigma
 =\int_{\R}P(c+iy)\overline{Q(c+iy)}\sigma(y)\,dy.
\]
The evaluation map is
$E_M:\mathcal P_M\longrightarrow\C^l$, $E_MP=(P(\alpha_j))_{j=0}^{l-1}$.
Its covariance $A_M=E_ME_M^*$ is in the stated Euclidean target metric;
no probability-mass rescaling is made.

\begin{lemma}[The original monic system and its mass]
The monic polynomials $p_n\in\C[S]$ obey
\[
 p_0=1,\quad p_1=S-c,\quad
 p_{n+1}=(S-c)p_n+n(n-\tfrac12)p_{n-1},\qquad
 \|p_n\|_\sigma^2=c_\sigma n!(\vartheta)_n.
\]
\end{lemma}
\begin{proof}
For $a>0$, set $f_a(t)=\exp(at-e^t)$. The substitution $r=e^t$ in
the Gamma integral gives
$\Gamma(a+iy/2)=\int_{\R}f_a(t)e^{iyt/2}\,dt$.
Plancherel with this explicit frequency $y/2$ gives, for real $u$,
\begin{align*}
 \int_{\R}\frac{|\Gamma(a+iy/2)|^2}{2\pi}e^{iuy}\,dy
 &=2\int_{\R}f_a(t)f_a(t+2u)\,dt\\
 &=\frac{2\Gamma(2a)e^{2au}}{(1+e^{2u})^{2a}}
 =2^{1-2a}\Gamma(2a)(\cosh u)^{-2a}.
\end{align*}
In particular the mass for $a=1/4$ is $c_\sigma=\sqrt{2\pi}$.
Every derivative of $f_a$ belongs to $L^2(\R)$, since it is $f_a$
times a polynomial in $e^t$; Plancherel therefore gives every required
polynomial moment of $|\Gamma(a+iy/2)|^2$.

Define, for real $z$ near zero,
\[
 F(x,z)=(1-z^2)^{-1/4}\exp(x\operatorname{arctanh}z)
 =\sum_{n\ge0}p_n(c+x)\frac{z^n}{n!}.
\]
The identity $(1-z^2)\partial_zF=(x+z/2)F$ gives the displayed recurrence
and monicity. For real $z,w$ near zero the Fourier formula gives
\begin{align*}
 \int_{\R}F(iy,z)\overline{F(iy,w)}\sigma(y)\,dy
 &=(1-z^2)^{-1/4}(1-w^2)^{-1/4}c_\sigma
   \bigl(\cosh(\operatorname{arctanh}z-
                       \operatorname{arctanh}w)\bigr)^{-1/2}\\
 &=c_\sigma(1-zw)^{-1/2}.
\end{align*}
Differentiating at $(z,w)=(0,0)$ is justified by the polynomial moments
just established. It gives orthogonality and the stated norms, including
the full factor $c_\sigma$.
\end{proof}

At $x=2j+1/2$, the same generating function is
$(1+z)^j(1-z)^{-j-1/2}$. Consequently
\begin{equation}\label{free:eq:eval}
 p_n(\alpha_j)=(\vartheta)_n Q_j(n),\qquad
 Q_j(n)=\sum_{r=0}^j\binom jr\frac{2^r n^{\underline r}}{(\vartheta)_r}.
\end{equation}
Indeed summing the right side divided by $n!$ gives
\[
 \sum_{r=0}^j\binom jr(2z)^r(1-z)^{-\vartheta-r}
 =(1+z)^j(1-z)^{-j-\vartheta}.
\]
Thus $Q_j$ has degree $j$ and leading coefficient $2^j/(\vartheta)_j$,
and the actual covariance is
\begin{equation}\label{free:eq:cov}
 (A_M)_{ij}=c_\sigma^{-1}\sum_{n=0}^M
          \frac{(\vartheta)_n}{n!}Q_i(n)Q_j(n).
\end{equation}

\section{Adjoint polynomial maps in both original measures}
For this section only, allow any $\vartheta>0$ and any integer $M\ge0$.
Put $w_n=(\vartheta)_n/n!$ and $K_M=(\vartheta)_{M+1}/M!$.
The discrete space uses $\sum_{n=0}^Mw_n\delta_n$ and the continuous
space on $[0,1]$ uses $K_Mt^{\vartheta-1}dt$.
The positive joint measure
\[
 K_M\binom Mn t^{n+\vartheta-1}(1-t)^{M-n}\,dt
 \quad (0\le n\le M)
\]
has the continuous marginal by the binomial theorem and the discrete
marginal because
\[
 K_M\binom Mn\mathrm B(n+\vartheta,M-n+1)
 =\frac{\Gamma(n+\vartheta)}{\Gamma(\vartheta)n!}=w_n.
\]
Its common mass is $K_M/\vartheta$; neither measure is replaced by a
probability measure. It gives the exact adjunction
\begin{align*}
 B:\C^{\{0,\ldots,M\}}&\longrightarrow L^2([0,1],K_Mt^{\vartheta-1}dt),
 & (BP)(t)&=\sum_{n=0}^M\binom Mn t^n(1-t)^{M-n}P(n),\\
 B^*(t^r)(n)&=\frac{(n+\vartheta)_r}{(M+\vartheta+1)_r},
 &\langle BP,R\rangle_{\rm cont}&=\langle P,B^*R\rangle_{\rm disc}.
\end{align*}
The second formula follows from the ratio of the two beta integrals,
so it proves both the map and its adjoint in the stated metrics.

\begin{lemma}[Exact discrete monic norms]
For $0\le j\le M$, the squared norm of the discrete monic orthogonal
polynomial of degree $j$ is
\begin{equation}\label{free:eq:norm}
 \nu_j(M)=\frac{(\vartheta)_{M+j+1}(j!)^2}
 {(M-j)!(2j+\vartheta)(j+\vartheta)_j^2}.
\end{equation}
\end{lemma}
\begin{proof}
The continuous monic orthogonal polynomial is
\[
 P_j(t)=\frac{(-1)^j}{(j+\vartheta)_j}
 t^{1-\vartheta}\frac{d^j}{dt^j}
       \{t^{j+\vartheta-1}(1-t)^j\}.
\]
Its leading coefficient is one. Integration by parts $j$ times proves
orthogonality against every smaller degree. For each derivative of
order $s\le j-1$ of the bracket, the factor at zero is
$O(t^{j+\vartheta-1-s})$, with exponent at least $\vartheta>0$;
at one the remaining vanishing order is at least $j-s\ge1$.
Every boundary term therefore vanishes. Applying the same integration
by parts to its own norm gives
\[
 \|P_j\|_{\rm cont}^2
 =\frac{K_Mj!}{(j+\vartheta)_j}
       \mathrm B(j+\vartheta,j+1)
 =\frac{K_M(j!)^2}{(2j+\vartheta)(j+\vartheta)_j^2}.
\]
Let $Q_j^{\rm mon}$ be the discrete monic polynomial. The exact identity
$B(n^{\underline j})=M^{\underline j}t^j$ and the displayed formula
for $B^*$ prove preservation of degrees and their leading coefficients.
Adjunction against every lower degree then proves
\[
 BQ_j^{\rm mon}=M^{\underline j}P_j,\qquad
 B^*P_j=\frac{Q_j^{\rm mon}}{(M+\vartheta+1)_j}.
\]
Evaluating $\langle BQ_j^{\rm mon},P_j\rangle$ on the two sides of the
adjunction yields
$\nu_j(M)=M^{\underline j}(M+\vartheta+1)_j\|P_j\|_{\rm cont}^2$.
Substitution gives \eqref{free:eq:norm}, also for $j=0$.
\end{proof}

Return to $\vartheta=1/2$. The triangular change from $n^j$ to $Q_j(n)$
in \eqref{free:eq:cov} has determinant $\prod_{j<l}2^j/(\vartheta)_j$;
the monic orthogonal change has determinant one. For every $M\ge l-1$,
\begin{equation}\label{free:eq:det}
 \det A_M=c_\sigma^{-l}\prod_{j=0}^{l-1}
 \frac{4^j(j!)^2(\vartheta)_{M+j+1}}
 {(\vartheta)_j^2(M-j)!(2j+\vartheta)(j+\vartheta)_j^2}.
\end{equation}
In particular all determinants are strictly positive, and
\begin{equation}\label{free:eq:inc}
 \frac{\det A_{M+1}}{\det A_M}
 =\prod_{j=0}^{l-1}\frac{M+j+\vartheta+1}{M+1-j}.
\end{equation}

\section{All four endpoints and a sharper finite expansion}
For every integer $q\ge1$ put
\[
 C_{q,l}=\log\frac{\det A_{2q+l-1}\det A_{2q+l}}
 {\det A_{q+l-1}\det A_{q+l}},\qquad B_l=l(l-\tfrac12).
\]
The constants in \eqref{free:eq:det} are retained until this exact signed
combination is taken. Set $a_j=l-j$, $b_j=l+j+1/2$, and
$f_j(x)=\log((x+b_j)/(x+a_j))$. Formula \eqref{free:eq:inc} gives
\begin{equation}\label{free:eq:windows}
 C_{q,l}=\sum_{j=0}^{l-1}\sum_{n=q}^{2q-1}
                         \{f_j(n)+f_j(n+1)\}.
\end{equation}
All summands are positive. Applying
$u/(1+u)\le\log(1+u)\le u$ to these exact two windows, whose upper
denominators are at least $n+1$ and lower denominators at most
$n+2l+1/2$, gives
\[
 2B_l\log\frac{2q+2l+1/2}{q+2l+1/2}
 \le C_{q,l}\le2B_l\log2.
\]
Here the decreasing-integrand comparisons are respectively
$\sum_{n=q}^{2q-1}(n+d)^{-1}\ge\int_q^{2q}(x+d)^{-1}dx$ and
$\sum_{n=q}^{2q-1}(n+1)^{-1}\le\int_q^{2q}x^{-1}dx$.
These justify the incoming finite bounds without an asymptotic step.

\begin{theorem}[Uniform finite signed expansion]\label{free:thm:expansion}
For every integer $q,l\ge1$,
\begin{equation}\label{free:eq:first}
 C_{q,l}=2B_l\log2-\frac{B_l(l+1/4)}q+R_{q,l},\qquad
 0\le R_{q,l}\le\frac{B_l(14l^2+5l+3)}{16q^2}.
\end{equation}
There is also the exact second-order enclosure
\begin{align}
 C_{q,l}
 &=2B_l\log2-\frac{B_l(l+1/4)}q
  +\frac{l^2(2l-1)(14l+5)}{32q^2}+E_{q,l},\label{free:eq:second}\\
 -\frac{7l(2l-1)(4l+1)(12l^2+2l-1)}{768q^3}
 &\le E_{q,l}\le\frac{3B_l}{16q^2}.\label{free:eq:error}
\end{align}
\end{theorem}
\begin{proof}
For a twice differentiable $f$, integration by parts on one unit interval
gives the exact trapezoidal identity
\[
 \frac{f(n)+f(n+1)}2-\int_n^{n+1}f(x)dx
 =\frac12\int_0^1s(1-s)f''(n+s)ds.
\]
For our function
$f_j(x)=\int_{a_j}^{b_j}(x+t)^{-1}dt$, its second derivative is positive.
Since $s(1-s)/2\le1/8$, summing the identity and using
\eqref{free:eq:windows} gives
\begin{align*}
 C_{q,l}&=2\sum_j\int_q^{2q}f_j(x)dx+T_{q,l},\\
 0\le T_{q,l}
 &\le\frac14\sum_j\{f_j'(2q)-f_j'(q)\}\\
 &=\frac14\sum_j\int_{a_j}^{b_j}
                  \left(\frac1{(q+t)^2}-\frac1{(2q+t)^2}\right)dt
 \le\frac{3B_l}{16q^2}.
\end{align*}
The final integrand decreases for $t\ge0$ and at zero equals $3/(4q^2)$;
also $\sum_j(b_j-a_j)=B_l$. Fubini, in positive finite integrals, gives
\[
 2\sum_j\int_q^{2q}f_j(x)dx
 =2\sum_j\int_{a_j}^{b_j}g_q(t)dt,
 \qquad g_q(t)=\log\frac{2q+t}{q+t}.
\]
For $t\ge0$, direct differentiation shows
\[
 g_q(0)=\log2,\quad g_q'(0)=-\frac1{2q},\quad
 0<g_q''(t)\le\frac3{4q^2},\quad
 -\frac7{4q^3}\le g_q'''(t)<0.
\]
Taylor's formula with integral remainder therefore gives both
\begin{gather*}
 0\le g_q(t)-\log2+\frac{t}{2q}\le\frac{3t^2}{8q^2},\\
 -\frac{7t^3}{24q^3}
 \le g_q(t)-\log2+\frac{t}{2q}-\frac{3t^2}{8q^2}\le0.
\end{gather*}
The finite power sums, obtained by expansion and
$\sum_{j<l}j=l(l-1)/2$, $\sum_{j<l}j^2=l(l-1)(2l-1)/6$,
$\sum_{j<l}j^3=l^2(l-1)^2/4$, are
\begin{align*}
 \sum_j\int_{a_j}^{b_j}t\,dt&=B_l(l+\tfrac14),\\
 \sum_j(b_j^3-a_j^3)&=\frac{l^2(2l-1)(14l+5)}8,\\
 \sum_j(b_j^4-a_j^4)&=
       \frac{l(2l-1)(4l+1)(12l^2+2l-1)}{16}.
\end{align*}
Integration of the first Taylor enclosure, together with the bound on
$T_{q,l}$, gives \eqref{free:eq:first}, because
\[
 \frac14\sum_j(b_j^3-a_j^3)+\frac{3B_l}{16}
 =\frac{B_l(14l^2+5l+3)}{16}.
\]
Integration of the second enclosure gives \eqref{free:eq:second} and
\eqref{free:eq:error}; its lower bound is
$-7\sum_j(b_j^4-a_j^4)/(48q^3)$ and its upper bound is the retained
positive trapezoidal error. This proves both finite formulas.
\end{proof}

\section{The stipulated packet degree is retained in the limits}
The programme uses precisely
\[
 q=(4l+2)^2\bigl(1+(4l+1)(m-1)\bigr),\qquad m\ge1.
\]
For $m=1$, substitution into \eqref{free:eq:second} gives
\begin{equation}\label{free:eq:simple}
 C_{q,l}=2l^2\log2-l(\log2+\tfrac1{16})
                 +\frac{167}{2048}+O(l^{-1}).
\end{equation}
Every error in this statement is bounded by \eqref{free:eq:error} and a
rational remainder from the displayed value of $q$. In detail,
\[
 \frac{B_l(l+1/4)}{(4l+2)^2}
 =\frac l{16}-\frac5{64}+O(l^{-1}),\qquad
 \frac{l^2(2l-1)(14l+5)}{32(4l+2)^4}
 =\frac7{2048}+O(l^{-1}),
\]
and \eqref{free:eq:error} is $O(l^{-1})$. Hence $C_{q,l}/q\to(\log2)/8$
and the constant term in \eqref{free:eq:simple} is $5/64+7/2048=167/2048$.

For each fixed integer $m>1$, $q/l^3\to64(m-1)$. Formula
\eqref{free:eq:first} has $R_{q,l}=O(l^{-2})$, while
$B_l(l+1/4)/q\to1/(64(m-1))$. Thus
\begin{equation}\label{free:eq:multiple}
 C_{q,l}-2l(l-\tfrac12)\log2\longrightarrow-\frac1{64(m-1)},
 \qquad
 \frac{lC_{q,l}}q\longrightarrow\frac{\log2}{32(m-1)}.
\end{equation}
The relation-dependent matrices at the same four endpoints remain in
the exact total return. Equations \eqref{free:eq:simple}--\eqref{free:eq:multiple}
evaluate its explicitly identified free summand, with a uniform finite
error already proved above; they impose no value on that remaining
summand.
